\chapter{Height fields: displacement, bump and normal maps}
\label{ch:height}

A grey image $h(u,v)\in[0,1]$ can describe the relief of a surface: light texels are high,
dark texels are low. Such an image is a \emph{height field}, and there are two very different
ways to use it. A \emph{displacement map} moves the surface; a \emph{bump map} leaves the
surface where it is and only changes the normal used for lighting, so that the surface
\emph{looks} as if it had been moved. The materials of Texture Explorer come with a height
map labelled ``displacement'', and the program uses the same image in both roles; switching
each role on and off separately is the quickest way to understand the difference
(Figure~\ref{fig:bumpvsdisp}).

\begin{figure}[ht]
\centering
\begin{tikzpicture}[scale=1.25,
  hfun/.style={domain=0:6,samples=120,smooth}]
  % ---- bump mapping
  \begin{scope}
    \draw[black!30,dashed,hfun] plot (\x,{0.2*sin(2*\x r)+0.08*sin(5*\x r)});
    \draw[very thick] (0,0) -- (6,0);
    \foreach \x in {0.3,0.75,...,5.9} {
      \pgfmathsetmacro{\d}{0.4*cos(2*\x r)+0.4*cos(5*\x r)}
      \pgfmathsetmacro{\nx}{-\d/sqrt(1+\d*\d)}
      \pgfmathsetmacro{\ny}{1/sqrt(1+\d*\d)}
      \draw[-{Stealth},npcolor,thick] (\x,0) -- ++(0.7*\nx,0.7*\ny);
    }
    \node[font=\small] at (3,-0.55) {\textbf{bump mapping}: the surface stays flat; only the normals tilt, following the slope of $h$ (dashed)};
  \end{scope}
  % ---- displacement mapping
  \begin{scope}[shift={(0,-2.3)}]
    \draw[very thick,hfun] plot (\x,{0.2*sin(2*\x r)+0.08*sin(5*\x r)});
    \draw[black!30,dashed] (0,0) -- (6,0);
    \foreach \x in {0.3,0.75,...,5.9} {
      \pgfmathsetmacro{\h}{0.2*sin(2*\x r)+0.08*sin(5*\x r)}
      \pgfmathsetmacro{\d}{0.4*cos(2*\x r)+0.4*cos(5*\x r)}
      \pgfmathsetmacro{\nx}{-\d/sqrt(1+\d*\d)}
      \pgfmathsetmacro{\ny}{1/sqrt(1+\d*\d)}
      \draw[-{Stealth},npcolor,thick] (\x,\h) -- ++(0.7*\nx,0.7*\ny);
    }
    \node[font=\small] at (3,-0.55) {\textbf{displacement mapping}: the surface itself moves; silhouette and occlusion change};
  \end{scope}
\end{tikzpicture}
\caption{The same height field (here a sum of two sines) used in its two roles, seen in
cross-section. With bump mapping the normals are exactly those of the displaced surface,
but they are attached to the undisplaced one.}
\label{fig:bumpvsdisp}
\end{figure}

\section{Displacement mapping}
\label{sec:displacement}

Let $\vect P(u,v)$ be a surface with unit normal $\Nv(u,v)$. Displacing it by the height
field gives the new surface
\begin{equation}
  \label{eq:displace}
  \vect P'(u,v) = \vect P(u,v) + s\,\bigl(h(u,v)-m\bigr)\,\Nv(u,v),
\end{equation}
where $s$ is the \emph{displacement scale}, a length, and $m$ the \emph{mid-level}, the height
that leaves the surface in place. With the program's default $m=\frac12$, dark texels push the
surface in and light texels push it out; with $m=0$ the surface only grows.

Displacement mapping was introduced by Cook \cite{cook1984}. On the GPU it is applied per vertex in the vertex shader (Section~\ref{sec:vs}).
This is the most direct technique and has two consequences worth remembering.
\begin{itemize}
\item \emph{The mesh must be dense.} A vertex samples the height field once; detail between
vertices is lost. Texture Explorer builds every patch as a grid of $256\times256$ cells,
about $66\,000$ vertices, and samples the height at the matching mipmap level
(Chapter~\ref{ch:sampling}). Production renderers use hardware tessellation to make the mesh
dense only where it is needed.
\item \emph{Normals are not updated.} Moving the vertices does not change the normals that the
mesh carries. A displaced surface shaded with its old normals has the right silhouette and the
wrong lighting. The fix is to compute the normal of the displaced surface, which is exactly
what bump mapping does---which is why displacement and bump (or normal) maps are almost always
used together.
\end{itemize}

\paragraph{Cracks.} Equation~\eqref{eq:displace} pushes each point along \emph{its} normal.
Where a solid has a crease---the edges of the pyramid, the rims of the cylinder---the two
patches that meet there have different normals, and vertices at the same position move in
different directions: the surface tears open. Texture Explorer leaves the cracks visible on
purpose. The usual remedies are to displace along an averaged normal at the crease, or to
model creases with a small bevel.

\section{Bump mapping}
\label{sec:bump}

In 1978 James Blinn observed that the appearance of small wrinkles is due almost entirely to
the way they tilt the normal, not to the way they move the surface \cite{blinn1978}. So we can
\emph{compute} the normal of the displaced surface \eqref{eq:displace} and use it to shade the
\emph{undisplaced} one.

\subsection{Blinn's derivation}
Differentiating \eqref{eq:displace} (and writing $\tilde h = s(h-m)$),
\[
  \vect P'_u = \vect P_u + \tilde h_u\,\Nv + \tilde h\,\Nv_u,\qquad
  \vect P'_v = \vect P_v + \tilde h_v\,\Nv + \tilde h\,\Nv_v .
\]
For small bumps $\tilde h$ itself is tiny while its derivatives $\tilde h_u$, $\tilde h_v$ may
not be (a small but steep groove), so the terms $\tilde h\,\Nv_u$ and $\tilde h\,\Nv_v$ are
dropped. The normal of the displaced surface is then parallel to
\[
  \vect P'_u\times\vect P'_v \approx
  \vect P_u\times\vect P_v + \tilde h_u\,(\Nv\times\vect P_v) + \tilde h_v\,(\vect P_u\times\Nv),
\]
(the term $\tilde h_u\tilde h_v\,\Nv\times\Nv$ vanishes), which is Blinn's perturbed normal.

\subsection{In an orthonormal frame}
Suppose that, at the point, $\vect P_u = a\,\Tv$ and $\vect P_v = b\,\Bv$ with $\Tv$ and $\Bv$
orthonormal and perpendicular to $\Nv$. Then $\tilde h_u/a$ is the rate of change of the height
per unit \emph{length} along $\Tv$, and similarly $\tilde h_v/b$ along $\Bv$. Substituting
and simplifying---the result is the same whether the frame is right- or left-handed---gives
\begin{equation}
  \label{eq:bumpframe}
  \vect n' \;\propto\; \Nv \;-\; \frac{\tilde h_u}{a}\,\Tv \;-\; \frac{\tilde h_v}{b}\,\Bv .
\end{equation}
A geometric argument reaches the same result without cross products, and without worrying
about handedness. In the frame $(\Tv,\Bv,\Nv)$ the displaced surface is locally the graph
$z = g(x,y)$, whose tangent vectors are $(1,0,g_x)$ and $(0,1,g_y)$. The vector
$(-g_x,-g_y,1)$ is perpendicular to both:
\[
  (-g_x,-g_y,1)\cdot(1,0,g_x) = -g_x + g_x = 0,\qquad
  (-g_x,-g_y,1)\cdot(0,1,g_y) = -g_y + g_y = 0 .
\]
So the normal leans \emph{against} the slope: where the height increases to the right, the
normal tilts to the left. On a flat region $g_x=g_y=0$ and the normal is $(0,0,1)$, the
unperturbed $\Nv$.

\subsection{The formula used in Texture Explorer}
\label{sec:bumpformula}
Texture Explorer measures the slope in \emph{texels}: $h_x$ is the change of the bump map per
texel in the $u$ direction and $h_y$ per texel in the $v$ direction. A single \emph{bump
strength} $k$, set with a slider, converts it into a slope of the surface. The tangent-space
normal is
\begin{equation}
  \label{eq:nts}
  \nts = \frac{(-k\,h_x,\ -k\,h_y,\ 1)}{\lVert(-k\,h_x,\ -k\,h_y,\ 1)\rVert}.
\end{equation}
Comparing with \eqref{eq:bumpframe}, $k\,h_x$ plays the role of $\tilde h_u/a$, so
$k = sW/a$ where $W$ is the width of the texture in texels: $k$ hides the displacement scale
\emph{and} the stretch $a=\lVert\vect P_u\rVert$ of the parametrization. Using one $k$ for the
whole solid is a deliberate simplification: the bumps look equally deep everywhere, even where
the texture is stretched, as on the pyramid near its base. A physically exact bump map would
divide by the local stretch.

\section{Estimating the slope}
\label{sec:centraldiff}

\begin{figure}[ht]
\centering
\begin{tikzpicture}[scale=1.1]
  \draw[black!25] (0,0) grid (3,3);
  \fill[black!8] (1,1) rectangle (2,2);
  \node at (1.5,1.5) {$h(x,y)$};
  \node[tcolor] at (0.5,1.5) {$-$};
  \node[tcolor] at (2.5,1.5) {$+$};
  \node[bcolor] at (1.5,2.5) {$-$};
  \node[bcolor] at (1.5,0.5) {$+$};
  \node[right,align=left,font=\small] at (3.3,1.5)
    {$h_x \approx \tfrac12\bigl(h(x{+}1,y)-h(x{-}1,y)\bigr)$\\[4pt]
     $h_y \approx \tfrac12\bigl(h(x,y{+}1)-h(x,y{-}1)\bigr)$\\[4pt]
     ($y$, like $v$, grows downwards)};
\end{tikzpicture}
\caption{The central-difference stencil. The texel itself does not take part.}
\label{fig:stencil}
\end{figure}

The derivatives of the bump map are estimated by \emph{central differences}, one texel to each
side (Figure~\ref{fig:stencil}). Taylor's theorem shows why they are preferable to the simpler
forward difference $h(x+1)-h(x)$:
\[
  \frac{h(x+1)-h(x-1)}{2} = h'(x) + \frac{h'''(x)}{6} + \dots,\qquad
  h(x+1)-h(x) = h'(x) + \frac{h''(x)}{2} + \dots
\]
The central difference is exact for quadratics and does not shift the normal by half a texel.
Both the shader and the CPU probe take the four samples with bilinear filtering at offsets of
exactly one texel, so they agree to rounding error.

\begin{tryit}
Pin the probe on the edge of a brick. Read the slopes $h_x$, $h_y$ and the tangent-space
normal in the Inspector and check equation \eqref{eq:nts} by hand: with $k=4$,
$h_x=0.0038$ and $h_y=0.0093$ give $\nts\approx(-0.015,-0.037,0.999)$. Then change the bump
strength and watch the cyan arrow $\vect n'$ in the 3D view lean further.
\end{tryit}

\section{Normal maps}
\label{sec:normalmaps}

Computing derivatives in the shader costs four texture reads per pixel. A \emph{normal map}
stores the result instead: each texel holds $\nts$, encoded as a colour by
\[
  \mathit{rgb} = \tfrac12\,\nts + \tfrac12,\qquad \nts = 2\,\mathit{rgb}-1 .
\]
Because the normals of a mostly flat surface are close to $(0,0,1)$, normal maps have the
characteristic lavender colour $(0.5,0.5,1)$. They are usually not computed from a height map
but \emph{baked} from a detailed (``high-poly'') model onto a simple one, and can represent
details that are not height fields at all, such as overhangs.

Two conventions exist, differing in the sign of the green channel (Section~\ref{sec:uvconv}):
\emph{DirectX} normal maps store the component along $+v$ with $v$ downwards, \emph{OpenGL}
ones along $+v$ with $v$ upwards. Using a map of the wrong convention makes the relief look
lit from below. Texture Explorer's tangent-space $y$ axis is $+v$ downwards, so it decodes the
DirectX variant directly; the Inspector shows it next to the normal computed from the bump
map. The numbers differ---the authors of the material baked their map with their own strength
and from their own high-resolution geometry---but the direction in which the normal leans
agrees, which is a good check of the conventions.

Normal maps have a subtle filtering problem: averaging unit vectors (as mipmapping does)
produces shorter vectors, and renormalizing them loses the information that the surface was
bumpy at that scale. Advanced techniques (Toksvig's, or LEAN mapping) turn that lost length into
extra roughness.

\section{Summary}

\begin{table}[ht]
\centering\small
\begin{tabularx}{\linewidth}{@{}l>{\raggedright}X>{\raggedright}X>{\raggedright\arraybackslash}X@{}}
\toprule
 & \textbf{Bump map} & \textbf{Normal map} & \textbf{Displacement map}\\
\midrule
Stores & height $h$ & tangent-space normal $\nts$ & height $h$\\
Changes the shading normal & yes, from the slope of $h$ & yes, directly & no (unless normals are recomputed)\\
Changes silhouette and occlusion & no & no & yes\\
Cost & 4 extra samples per pixel & 1 sample per pixel & dense mesh or tessellation\\
Typical artefacts & flat silhouettes; light leaking on the dark side & same, plus wrong green convention & cracks at creases; swimming if under-sampled\\
In Texture Explorer & \code{BumpNormalTS} in the pixel shader & shown for comparison only & \code{VSMain} in the vertex shader\\
\bottomrule
\end{tabularx}
\caption{Three ways to use relief in a material.}
\label{tab:reliefs}
\end{table}

Bump and normal maps share one more artefact. A perturbed normal may face the light at points
where the real surface faces away from it, so the dark side of the solid sparkles. Texture
Explorer attenuates the lighting with a steep ramp of the \emph{geometric} cosine
$\Nv\cdot\vect L$ (Section~\ref{sec:horizon}), a crude form of self-shadowing.
